\section{Introduction}
Finite descriptions of an experiment can be compared only after the
questions to be answered and the permitted execution have been fixed.
A two-valued checkpoint target may require exponentially many causal
states when the relevant question arrives after the observation has been
discarded. A finite statistic may preserve every bounded decision problem
approximately without producing an accurate realization of a whole
posterior measure. These are not competing definitions of the same
error. They are different slices of one constrained decision problem.

We attach to a causal experiment its \emph{marked decision spectrum}:
the attainable risk vectors, indexed by the retained resources and by
families of decision tasks. One common encoder is chosen before the
task, whereas its decoder may depend on the task. The spectrum retains
this order of quantifiers. It separates persistent state, physical output,
acquisition length, program description, numerical precision and transient
work. Under a resource-controlled causal equivalence it is invariant,
up to the certified budget transformations. For a compact dominated
terminal experiment, its decision-complete slice is exactly the optimal
reverse deficiency of a common finite-output encoding. Its quadratic
Bayesian slice is a compatible-partition distortion problem.

\begin{theorem}[Representation and realization of the decision spectrum]
\label{thm:v8-main}
The spectrum defined in Section~\ref{sec:v8-spectrum} has the following
three properties.
\begin{enumerate}
\item Initialized, parameter-independent causal simulations transport
its attainable risk sets with their full resource transformers. Mutually
exact finite-state simulations preserve every existing positive
cardinality exponent.
\item For a compact total-variation-continuous terminal family,
\[
 \Theta_E(S;\mathcal D_{\rm all})
          =\inf_{Q:X\to[S]}\delta(QE,E).
\]
A posterior coordinate formed with any positive dense reference prior
is sufficient for the entire parameter family and leaves this terminal
spectrum unchanged. Its causal analogue requires likelihood saturation
and recursive predictive closure; a finite quotient then supplies a
finite-state equivalence with an explicitly charged budget dilation.
\item For finite report alphabets and a fixed finite-horizon Bayesian
quadratic problem, the optimal causal excess is exactly the minimum
conditional-variance functional over successor-compatible partitions
of the history levels, with width at most $S$.
\end{enumerate}
The precise hypotheses and proofs are Theorems~\ref{thm:v8-spectrum},
\ref{thm:v8-duality}, \ref{thm:v8-sufficient},
\ref{thm:v8-causal-sufficiency} and \ref{thm:v8-congruence}.
\end{theorem}

The third assertion supplies a necessity theorem rather than another
sufficient coding scheme. In Example~\ref{ex:v8-index}, a random
$d$-bit word is followed by a random coordinate query. The checkpoint
mean has two values, but zero-error causal realization needs $2^d$
states. Its positive error is bounded below by the binary entropy
inverse. Thus static predictive geometry alone cannot determine a
causal budget. Compatibility of the partitions is the missing datum.

We then prove three quantitative consequences. First,
Theorem~\ref{thm:v8-certified} replaces exact real-energy ordering in a
predictive tree by a rational threshold construction. Taking the minimum
of certified energy bounds over every contiguous factor gives exact
suffix closure despite ties and numerical uncertainty. Its finite tree
has optimal distortion order, $O(\log S)$ acquisition depth and an
$O(S\log S+Sd p)$-bit description for a $d$-dimensional readout. The
acquisition is a finite prefix of a snapshot, not a free infinite word.

Second, Theorem~\ref{thm:v8-measureentropy} and
Corollary~\ref{cor:v8-nonlinear} give a matching $1/\log S$ presentation
law for the entire space of probability measures on an interval and for
a genuine nonlinear posterior class with arbitrary input prior at each
fixed time. A continuous finite-dimensional exact coordinate is
impossible on this class. Yet a single marginal Bernoulli probe admits
an $O(1/S)$ simulation. The separation proves that forgetting the output
mark can change the rate, even when the latent measure is unchanged.

Third, Theorem~\ref{thm:v8-seqgaussian} gives the decision-complete
terminal law $S^{-2/(nr)}$ for $n$ sequential $r$-dimensional Gaussian
innovations when only the current register is retained. The register
must hold the accumulated finite experiment; previous labels cannot be
stored in a free transcript. The positive Gaussian reconstruction and
the physical Poisson count-to-contact comparison are proved in full in
the supporting sections. The latter gives a direct downstream instance
with matched label budgets in a joint sample-size window.

\subsection{Relation to previous theory}
Predictive equivalence and minimal predictive states are classical
ideas; Shalizi and Crutchfield \cite{ShaliziCrutchfield2001} establish
optimality and uniqueness properties of causal-state representations.
The continuation congruence used here has the classical finite-automaton
antecedent \cite{Nerode1958}. We use neither of these observations as a
claim of a new exact minimal-state principle. The extra constraint is a
charged recursive presentation, with a fixed common encoder and a
quantitative loss. The finite-history variance formula identifies its
feasible partitions exactly.

The comparison of experiments through all bounded decision problems is
likewise the theory of Blackwell and Le Cam
\cite{Blackwell1953,LeCam1964,Torgersen1991}. Norberg
\cite{Norberg2002} treats comparison with filtered probability spaces.
Our terminal duality is proved to fix its compactness assumptions and
its encoder quantifiers, not to reassign priority to that criterion.
The new organization retains resource transformers as part of a causal
morphism and distinguishes saturated statistical quotients from
prior-relative predictive quotients.

Zero-delay coding of Markov sources has a substantial structural
literature. Linder--Y\"uksel \cite{LinderYuksel2014} and
Wood--Linder--Y\"uksel \cite{WoodLinderYuksel2017} study causal coding
policies, including stationary optimal policies and finite-memory
approximations. Their communication alphabet is not automatically the
cardinality of all retained information. In the latter paper, Theorems
3--4 of the author preprint give stationary average-cost optimality and
periodic finite-memory approximations. The sequential Gaussian result
here concerns a different, explicitly total-register terminal problem.

Functional quantization and separated-cylinder quantization provide the
static geometry \cite{LuschgyPages2004,Roychowdhury2014,KZ2017}; greedy
variable-to-fixed trees have the classical antecedent
\cite{Tunstall1967}. We retain those attributions. The robust factor
envelope proves that a near-optimal finite tree can be compiled without
an oracle for exact comparisons, while its suffix closure makes both
causal updates executable. Nonlinear-filter contraction and finite-model
approximation have extensive precedents, including
\cite{DKYAC2024}. Our additional lower bound concerns the full posterior
presentation task, not a new proof of average-cost control optimality.
The topological and entropy arguments also distinguish continuous
finite-dimensional coordinates from arbitrary Borel codings.

\subsection{Organization and quantitative conventions}
Sections~\ref{sec:v8-spectrum}--\ref{sec:v8-congruence} establish the
common spectrum, its sufficient-statistic representation and its causal
compatibility obstruction. Sections~\ref{sec:v8-certified}--\ref{sec:v8-gaussian}
prove the quantitative consequences. The remaining sections supply the
complete predictive-tree, filtering, Gaussian and count-experiment
proofs used by these consequences. They are mathematical dependencies,
not substitutes for a proof of the new representation statements.

The notation $a_S\asymp b_S$ means comparison by strictly positive
constants independent of $S$ in the stated parameter regime. Constants
for the tree class are uniform only when its separation, child-energy,
unary-chain and distortion constants have common bounds. Horizon,
covariance conditioning and rank dependencies are retained. Explicit
Poisson normal-approximation constants in the supporting proof are
nonoptimized. The physical A2 consequence is order-optimal in the label
budget inside its displayed joint window, not sample-optimal over all
observation protocols. The paper's title names the continuing programme;
the invariant proved here is the typed marked spectrum, not a claim that
unrelated spectral, kinetic or phase theorems are automatic corollaries.
